\documentclass[11pt,a4paper]{article}
\usepackage[utf8]{inputenc}
\usepackage[T1]{fontenc}
\IfFileExists{lmodern.sty}{\usepackage{lmodern}}{\IfFileExists{mathptmx.sty}{\usepackage{mathptmx}}{}}
\usepackage[ngerman]{babel}
\usepackage[a4paper,margin=27mm]{geometry}
\usepackage{amsmath,amssymb}
\usepackage{booktabs}
\usepackage{array}
\renewcommand{\floatpagefraction}{0.85}
\renewcommand{\textfraction}{0.1}
\IfFileExists{enumitem.sty}{\usepackage{enumitem}\setlist[itemize]{leftmargin=1.5em,itemsep=2pt}}{}
\usepackage{tikz}
\usetikzlibrary{arrows.meta,positioning,calc,shapes.misc}
\IfFileExists{microtype.sty}{\usepackage{microtype}}{}
\usepackage[hidelinks]{hyperref}
\setlength{\parskip}{0.35em}
\setlength{\parindent}{1.2em}
\newcommand{\SO}{Subjekt-Objekt}
\newcommand{\Werke}[1]{(W~2, S.~#1)}

\title{Subjekt-Objekt und die Identität der Identität und der Nichtidentität\\[0.4em]
\large Die Zusammensetzung des Absoluten, des Subjektiven und des Objektiven\\ in Hegels Differenzschrift (1801)}
\author{}
\date{}

\begin{document}
\maketitle

\section{Ort und Anlage}

Die Differenzschrift gliedert sich in eine Vorerinnerung, das Kapitel über die mancherlei Formen des gegenwärtigen Philosophierens, die Darstellung des Fichteschen Systems, die Vergleichung des Schellingschen Prinzips mit dem Fichteschen und den Anhang über Reinhold. Zwei Sachverhalte werden im Folgenden aus diesem Aufbau herausgehoben und am Wortlaut verfolgt.

Der erste Sachverhalt ist die Zusammensetzung. Das Absolute ist \SO{}; das Subjekt ist \SO{}; das Objekt ist \SO{}. Dieser Gedanke wird in der Vergleichung entwickelt \Werke{94--101} und bis zum Indifferenzpunkt beider Wissenschaften geführt \Werke{111~ff.}. Vorbereitet ist er durch das Urteil der Vorerinnerung, das Prinzip Fichtes erweise sich als ein bloß subjektives \SO{} \Werke{10~f.}, und durch die Darstellung des Fichteschen Systems.

Der zweite Sachverhalt ist die Formel: Das Absolute ist die Identität der Identität und der Nichtidentität. Sie steht ebenfalls in der Vergleichung \Werke{96}, aber ihre Bausteine stammen aus dem Kapitel über das Prinzip einer Philosophie in der Form eines absoluten Grundsatzes \Werke{35--41}: die beiden Sätze $A = A$ und $A = B$, ihre wechselseitige Bedingtheit und ihre Antinomie.

Die Vermutung, die Zusammensetzung mache den zweiten Teil der Schrift aus, trifft in dieser Form zu: Die Vergleichung ist der Ort, an dem die doppelte Setzung des \SO{} verlangt, begründet und zu den beiden Wissenschaften ausgefaltet wird. Beide Sachverhalte hängen dort so eng zusammen, dass die Formel als die Satzform dessen gelesen werden kann, was die Zusammensetzung als Struktur sagt. Dies wird im letzten Abschnitt ausgeführt.

Zitiert wird nach Band~2 der Theorie-Werkausgabe (Suhrkamp 1986), abgekürzt W~2. Die Seitenzahlen folgen der dort eingedruckten Paginierung.

\section{Erster Sachverhalt: Das Absolute, das Subjekt, das Objekt als \SO{}}

\subsection{Der Ausgangspunkt: Entzweiung}

Der Ausgangspunkt ist nicht ein Prinzip, sondern ein Riss. \glqq Entzweiung ist der Quell des Bedürfnisses der Philosophie\grqq{} \Werke{20}. Die Bildung des Zeitalters hat das, was Erscheinung des Absoluten ist, vom Absoluten isoliert und als Selbständiges fixiert. Die feste Gestalt dieser Entzweiung ist die Entgegensetzung von Subjekt und Objekt: von Ich und Welt, Intelligenz und Natur, Freiheit und Notwendigkeit. Die Aufgabe der Philosophie ist, die Entzweiung aufzuheben, das Absolute für das Bewusstsein zu konstruieren. Ob das gelingt, entscheidet sich daran, wie die beiden Entgegengesetzten in das Absolute gesetzt werden.

\subsection{Das Urteil über Fichte: ein subjektives \SO{}}

Die Vorerinnerung fällt das Urteil, das dann die ganze Darstellung des Fichteschen Systems einlöst. Das Prinzip Fichtes ist die transzendentale Anschauung, Ich~=~Ich: das Selbstbewusstsein, das sich selbst zum Gegenstand hat. In ihm ist das Subjekt, das setzt, und das Objekt, das gesetzt wird, dasselbe. Das ist der Sinn des Wortes \SO{}: eine Identität, in der Subjekt und Objekt nicht nachträglich zusammengebracht, sondern von Anfang an eins sind. Das Prinzip ist damit spekulativ.

Aber sobald die Spekulation \glqq sich zum System bildet\grqq{}, verlässt sie ihr Prinzip. Das Deduzierte erhält die Form einer Bedingung des reinen Bewusstseins; Ich~=~Ich wird durch eine objektive Unendlichkeit bedingt, einen Zeit-Progress ins Unendliche, in dem sich die transzendentale Anschauung verliert. \glqq Das Prinzip, das \SO{} erweist sich als ein subjektives \SO{}\grqq{} \Werke{10~f.}. Das Wort \emph{subjektiv} sagt hier: Die Identität von Subjekt und Objekt ist selbst nur unter der Form des Subjekts gesetzt. Sie ist das Ich, und ihr steht das Objekt als das bloß Gesetzte gegenüber.

Die Darstellung führt das aus. \glqq Weil Ich subjektives \SO{} ist, so bleibt ihm eine Seite, von welcher ihm ein Objekt absolut entgegengesetzt ist, von welcher es durch dasselbe bedingt ist\grqq{} \Werke{68--73}. Die Natur wird nur als Nicht-Ich deduziert, als Schranke der freien Tätigkeit. Vom Standpunkt des Ich erscheint das Objekt \glqq als ein nicht Sich-selbst-Bestimmendes, als ein absolut Bestimmtes\grqq{} \Werke{96}. Und darum verwandelt sich der Grundsatz: Ich~=~Ich wird zu \glqq Ich soll gleich Ich sein\grqq{} \Werke{97}. Die Identität ist nicht, sie wird gefordert. Das Absolute des Systems ist ein Sollen.

\subsection{Die Forderung der Vergleichung}

Die Vergleichung setzt mit der Diagnose ein: Bei Fichte tritt Subjekt~=~Objekt aus der Identität heraus und vermag sich zu ihr nicht wiederherzustellen, weil das Differente ins Kausalitätsverhältnis versetzt wurde. Das Prinzip der Identität wird nicht Prinzip des Systems. Dann folgt die Forderung, die den ersten Sachverhalt eröffnet:

\begin{quote}
\glqq Daß absolute Identität das Prinzip eines ganzen Systems sei, dazu ist notwendig, daß das Subjekt und Objekt beide als \SO{} gesetzt werden. Die Identität hat sich im Fichteschen System nur zu einem subjektiven \SO{} konstituiert. Dies bedarf zu seiner Ergänzung eines objektiven \SO{}s, so daß das Absolute sich in jedem der beiden darstellt, vollständig sich nur in beiden zusammen findet, als höchste Synthese in der Vernichtung beider, insofern sie entgegengesetzt sind, als ihr absoluter Indifferenzpunkt beide in sich schließt, beide gebiert und sich aus beiden gebiert.\grqq{} \Werke{94}
\end{quote}

Hier steht die ganze Struktur auf einmal. Zu lesen ist sie in drei Schritten.

Erstens: Subjekt und Objekt sind beide \SO{}. Was Fichte nur vom Ich sagt, gilt vom Objekt ebenso. Die Natur ist nicht bloß Objekt, sondern selbst Identität von Subjekt und Objekt in der Form des Objekts.

Zweitens: Das Absolute stellt sich in jedem der beiden dar und findet sich vollständig nur in beiden zusammen. Jede Seite ist das Ganze, aber unter einer Form; die Vollständigkeit des Ganzen ist keine der beiden Formen, sondern ihr Zusammen.

Drittens: Das Absolute ist Indifferenzpunkt, der beide gebiert und sich aus beiden gebiert. Die Richtung der Erzeugung ist doppelt. Das Ganze bringt die Seiten hervor, und die Seiten bringen das Ganze hervor.

\subsection{Gleiches Recht, gleiche Notwendigkeit, Teilung ins Unendliche}

Zwei Seiten später wird die Forderung begründet. Die Philosophie muss dem Trennen in Subjekt und Objekt sein Recht widerfahren lassen; sie kann die Getrennten aber nicht setzen, ohne sie im Absoluten zu setzen, denn sonst wären es \glqq rein Entgegengesetzte, die keinen anderen Charakter haben, als daß das eine nicht ist, insofern das andere ist\grqq{} \Werke{96}. Beide müssen im Absoluten gesetzt werden, und zwar mit gleichem Recht und gleicher Notwendigkeit, denn würde nur eines auf das Absolute bezogen, so wäre ihr Wesen ungleich gesetzt und die Vereinigung unmöglich. Fichte hat nur eines gesetzt, das Selbstbewusstsein, weil nur dieses sich selbst setzt. Dass diese Bevorzugung zufällig ist, folgt daraus, dass das Selbstbewusstsein selbst ein Bedingtes ist: Es ist die Vernunft in einer beschränkten Form.

Dann die Formulierung der Zusammensetzung:

\begin{quote}
\glqq Es müssen daher beide in das Absolute oder das Absolute in beiden Formen gesetzt werden und zugleich beide als Getrennte bestehen; das Subjekt ist hiermit subjektives \SO{}, das Objekt objektives \SO{}. Und weil nunmehr, da eine Zweiheit gesetzt ist, jedes der Entgegengesetzten ein sich selbst Entgegengesetztes ist und die Teilung ins Unendliche geht, so ist jeder Teil des Subjekts und jeder Teil des Objekts selbst im Absoluten, eine Identität des Subjekts und Objekts, -- jedes Erkennen eine Wahrheit, so wie jeder Staub eine Organisation. Nur indem das Objekt selbst ein \SO{} ist, ist Ich~=~Ich das Absolute.\grqq{} \Werke{96~f.}
\end{quote}

Damit ist die Rekursion ausgesprochen. Die Zweiheit von Subjekt und Objekt kehrt in jedem der beiden wieder, und sie kehrt in jedem Teil jedes der beiden wieder. Die Teilung geht ins Unendliche, ohne dass die Identität je verloren ginge: Jeder Teil ist im Absoluten und ist selbst eine Identität von Subjekt und Objekt. Der Satz vom Staub und von der Organisation ist die Probe: Das Kleinste ist nicht Rest der Teilung, sondern wieder das Ganze.

\begin{figure}[!htb]
\centering
\resizebox{\linewidth}{!}{%
\begin{tikzpicture}[
  every node/.style={align=center, font=\small},
  box/.style={draw, rounded corners=2pt, inner sep=5pt, minimum width=3.2cm},
  leaf/.style={draw, rounded corners=2pt, inner sep=3pt, font=\footnotesize},
  edge/.style={-{Latex[length=2mm]}, thick},
  back/.style={-{Latex[length=2mm]}, thick, dashed, gray!70!black},
  node distance=9mm and 7mm]
\node[box, fill=gray!12] (A) {\textbf{Absolutes}\\ \SO{}\\ Indifferenzpunkt: Identität \emph{und} Totalität};
\node[box, below left=14mm and 6mm of A] (S) {\textbf{Subjekt}\\ subjektives \SO{}\\ Ich = Ich, Intelligenz\\ \footnotesize Subjekt Substanz, Objekt Akzidens};
\node[box, below right=14mm and 6mm of A] (O) {\textbf{Objekt}\\ objektives \SO{}\\ Natur\\ \footnotesize Objekt Substanz, Subjekt Akzidens};
\node[leaf, below left=9mm and -6mm of S] (SS) {Subjekt\\ (im Subjekt)};
\node[leaf, below right=9mm and -6mm of S] (SO) {Objekt\\ (im Subjekt)};
\node[leaf, below left=9mm and -6mm of O] (OS) {Subjekt\\ (im Objekt)};
\node[leaf, below right=9mm and -6mm of O] (OO) {Objekt\\ (im Objekt)};
\draw[edge] (A) -- (S); \draw[edge] (A) -- (O);
\draw[edge] (S) -- (SS); \draw[edge] (S) -- (SO);
\draw[edge] (O) -- (OS); \draw[edge] (O) -- (OO);
\foreach \n in {SS,SO,OS,OO} {
  \draw[dotted, thick] (\n.south) -- ++(0,-6mm);
}
\node[below=7mm of $(SO.south)!0.5!(OS.south)$, font=\footnotesize\itshape] (inf) {Teilung ins Unendliche: jeder Teil wieder \SO{}, \glqq jeder Staub eine Organisation\grqq};
\draw[back] (S.north) to[bend left=25] node[left, font=\scriptsize, xshift=-1mm] {gebiert sich\\ aus beiden} (A.west);
\draw[back] (O.north) to[bend right=25] node[right, font=\scriptsize, xshift=1mm] {gebiert sich\\ aus beiden} (A.east);
\node[above=1mm of A, font=\scriptsize] {};
\end{tikzpicture}}
\caption{Die Zusammensetzung nach W~2, S.~94 und 96~f. Abwärts: Das Absolute \glqq gebiert beide\grqq{}, jede Seite ist wieder \SO{}, und die Teilung geht ins Unendliche. Aufwärts (gestrichelt): Das Absolute \glqq gebiert sich aus beiden\grqq{}. Die Unterschrift der Seiten nennt das Substantialitätsverhältnis nach S.~101.}
\label{fig:baum}
\end{figure}

\subsection{Reelle Entgegensetzung und die beiden Wissenschaften}

Warum muss die Zweiheit so tief hinabreichen? Weil sonst die Entgegensetzung nicht real wäre. \glqq Indem das Subjekt sowohl als das Objekt ein \SO{} sind, ist die Entgegensetzung des Subjekts und Objekts eine reelle Entgegensetzung; denn beide sind im Absoluten gesetzt und haben dadurch Realität\grqq{} \Werke{97}. Eine ideelle Entgegensetzung ist Werk der Reflexion, die von der absoluten Identität abstrahiert; eine reelle Entgegensetzung ist Werk der Vernunft, die beide Seiten \glqq nicht bloß in der Form des Erkennens, sondern auch in der Form des Seins\grqq{} setzt \Werke{98}. Die Pointe ist umgekehrt, als der Verstand erwartet: Gerade dadurch, dass beide Seiten dieselbe Identität in sich tragen, wird ihr Gegensatz wirklich. Ein Objekt, das bloß Objekt wäre, könnte dem Subjekt nur ideell, als sein Nicht, entgegenstehen. Ein Objekt, das selbst \SO{} ist, steht ihm als Wirkliches gegenüber.

\glqq Hierin besteht allein die wahre Identität, daß beide ein \SO{} sind, und zugleich die wahre Entgegensetzung, deren sie fähig sind\grqq{} \Werke{99}. Wo nicht beide \SO{} sind, bleibt die Entgegensetzung ideell und das Prinzip der Identität formal. Das System, dessen Prinzip formal ist, ist eine formelle Philosophie; wenn hingegen \glqq die Materie, das Objekt, selbst ein \SO{} ist, so kann die Trennung der Form und Materie wegfallen\grqq{}, und das Prinzip ist \glqq formales und materiales zugleich\grqq{} \Werke{99}.

Daraus folgen die beiden Wissenschaften. Weil Subjekt und Objekt beide als \SO{} gesetzt sind, ist \glqq jedes für sich nunmehr fähig, der Gegenstand einer besonderen Wissenschaft zu sein\grqq{} \Werke{99}. Die Transzendentalphilosophie ist die Wissenschaft vom subjektiven \SO{}, die Naturphilosophie die vom objektiven \SO{} \Werke{100~f.}. Jede abstrahiert vom Prinzip der anderen, aber diese Abstraktion ist keine Einseitigkeit, sondern Reinheit: Es wird nicht vom inneren Prinzip, sondern nur von der eigentümlichen Form der anderen Wissenschaft abstrahiert \Werke{100}. Natur und Selbstbewusstsein sind an sich so, wie sie in ihrer Wissenschaft gesetzt werden, \glqq weil es die Vernunft ist, die sie setzt, und die Vernunft setzt sie als \SO{}, also als das Absolute\grqq{}. Sie setzt sie so, \glqq weil sie es selbst ist, die sich als Natur und als Intelligenz produziert und sich in ihnen erkennt\grqq{} \Werke{100}.

Der Unterschied der beiden Seiten wird dann ganz bestimmt benannt. In beiden Wissenschaften stehen das Subjektive und das Objektive im Substantialitätsverhältnis; in der Transzendentalphilosophie ist die Intelligenz die absolute Substanz und die Natur Akzidens, in der Naturphilosophie die Natur die Substanz und die Intelligenz Akzidens \Werke{101}. Der Unterschied von Subjekt und Objekt ist also kein Unterschied des Gehalts, denn beide sind \SO{}, sondern ein Unterschied dessen, welche Seite in ihnen Substanz ist.

\subsection{Der Indifferenzpunkt}

Beide Wissenschaften sind Totalitäten, aber relative; als solche \glqq streben sie nach dem Indifferenzpunkt; als Identität und als relative Totalität liegt er überall in ihnen selbst, als absolute Totalität außer ihnen\grqq{} \Werke{111}. Da ihre Entgegensetzung reell ist, \glqq hängen sie als Pole der Indifferenz in dieser selbst zusammen; sie selbst sind die Linien, welche den Pol mit dem Mittelpunkt verknüpfen\grqq{}. Und: \glqq dieser Mittelpunkt ist selbst ein gedoppelter, einmal Identität, das andere Mal Totalität\grqq{} \Werke{111}. Die beiden Wissenschaften erscheinen daher als der Fortgang der Selbstkonstruktion der Identität zur Totalität.

\begin{figure}[!htb]
\centering
\resizebox{\linewidth}{!}{%
\begin{tikzpicture}[
  every node/.style={align=center, font=\small},
  pole/.style={draw, rounded corners=2pt, inner sep=5pt, minimum width=3.4cm},
  edge/.style={thick}]
\node[pole] (T) at (-6.9,0) {\textbf{Transzendental-}\\ \textbf{philosophie}\\ subjektives \SO{}\\ \footnotesize Intelligenz Substanz,\\ \footnotesize Natur Akzidens};
\node[pole] (N) at (6.9,0) {\textbf{Natur-}\\ \textbf{philosophie}\\ objektives \SO{}\\ \footnotesize Natur Substanz,\\ \footnotesize Intelligenz Akzidens};
\node[draw, rounded corners=8pt, fill=gray!12, inner sep=6pt] (I) at (0,0) {\textbf{Indifferenzpunkt}\\ \footnotesize gedoppelt:\\ \footnotesize Identität / Totalität};
\draw[edge] (T) -- node[above=0.5mm, font=\scriptsize] {die Wissenschaft selbst\\ ist die Linie} node[below=0.5mm, font=\scriptsize] {Pol der Indifferenz} (I);
\draw[edge] (I) -- node[above=0.5mm, font=\scriptsize] {die Wissenschaft selbst\\ ist die Linie} node[below=0.5mm, font=\scriptsize] {Pol der Indifferenz} (N);
\node[below=3mm of T, font=\scriptsize] {relative Totalität,\\ strebt nach dem Indifferenzpunkt};
\node[below=3mm of N, font=\scriptsize] {relative Totalität,\\ strebt nach dem Indifferenzpunkt};
\node[below=6mm of I, font=\scriptsize\itshape] {als Identität in beiden Polen,\\ als absolute Totalität außer ihnen};
\node[above=8mm of T, font=\scriptsize, draw, dashed, inner sep=3pt] (F) {Fichte: nur dieser Pol\\ als das Absolute gesetzt};
\draw[dashed, gray] (F) -- (T);
\end{tikzpicture}}
\caption{Die beiden Wissenschaften als Pole der Indifferenz nach W~2, S.~100~f. und 111. Der Mittelpunkt ist gedoppelt: Als Identität liegt er in jedem Pol, als Totalität außer beiden. Fichtes System setzt allein den linken Pol als absolut.}
\label{fig:pole}
\end{figure}

\subsection{Die Formeln, kommentiert}

Der Sachverhalt lässt sich in drei Formeln und einem Zusatz festhalten. Sie sind nicht Hegels Schreibweise, sondern die Verdichtung seiner Sätze.

\begin{center}
\renewcommand{\arraystretch}{1.35}
\begin{tabular}{@{}l >{\raggedright\arraybackslash}p{5.6cm} >{\raggedright\arraybackslash}p{7.4cm}@{}}
\toprule
 & Formel & Wortlaut, auf den sie sich stützt \\
\midrule
(A) & Absolutes $=$ \SO{} & \glqq als ihr absoluter Indifferenzpunkt beide in sich schließt\grqq{} (S.~94) \\
(S) & Subjekt $=$ (\SO{})$_{\text{Form: Subjekt}}$ & \glqq das Subjekt ist hiermit subjektives \SO{}\grqq{} (S.~96) \\
(O) & Objekt $=$ (\SO{})$_{\text{Form: Objekt}}$ & \glqq das Objekt objektives \SO{}\grqq{} (S.~96) \\
(R) & jeder Teil von S, jeder Teil von O $=$ \SO{} & \glqq die Teilung ins Unendliche geht\grqq{} (S.~96~f.) \\
\bottomrule
\end{tabular}
\end{center}

\paragraph{Zu (A): Das Absolute als \SO{}.}
Das Absolute ist nicht ein Drittes neben Subjekt und Objekt, sondern ihre Identität. Aber diese Identität ist nicht die Fichtesche, die nur unter der Form des Subjekts gesetzt ist. Sie ist der Indifferenzpunkt: die Stelle, an der die Differenz von Subjekt und Objekt gleichgültig wird, weil beide dasselbe sind. Dass der Mittelpunkt \glqq gedoppelt\grqq{} ist, Identität und Totalität, hält den Unterschied fest zwischen der Identität, die in jeder Seite liegt, und dem Ganzen, das nur beide zusammen sind. Der Indifferenzpunkt ist also nicht leer; er ist das Ganze, das sich in beide Formen auslegt.

\paragraph{Zu (S) und (O): Der Formindex.}
Die beiden Formeln unterscheiden sich nur im Index. Ihr Gehalt ist derselbe: \SO{}. Der Unterschied ist der der Form, und die Form wird von Hegel durch das Substantialitätsverhältnis bestimmt: welche der beiden Seiten in der Identität Substanz ist und welche Akzidens. Das ist die Grundlage dafür, dass Hegel die reelle Entgegensetzung \glqq nur quantitativ\grqq{} nennt \Werke{99}. Was zwischen Subjekt und Objekt differiert, ist nicht das Was, sondern das Überwiegen. Und deshalb ist das Absolute, das sich aus der quantitativen Differenz rekonstruiert, \glqq kein Quantum, sondern Totalität\grqq{}. Wäre der Unterschied qualitativ, dann könnte nur eine der beiden Seiten das Absolute sein, und man stünde wieder bei Fichte.

\paragraph{Zu (R): Die Rekursion.}
Der Zusatz macht aus den drei Formeln eine Struktur. Weil die Zweiheit gesetzt ist, ist \glqq jedes der Entgegengesetzten ein sich selbst Entgegengesetztes\grqq{}: Im Subjekt kehren Subjekt und Objekt wieder, im Objekt ebenso, und in jedem Teil davon wieder. Das ist eine Selbstähnlichkeit ohne Abbruch. Entscheidend ist, dass diese Unendlichkeit nicht die schlechte ist, die Hegel an Fichte tadelt. Dort verliert sich Ich~=~Ich in einem Progress, in dem die Identität nie erreicht wird \Werke{10~f.}. Hier ist die Identität auf jeder Stufe schon da: \glqq jeder Teil des Subjekts und jeder Teil des Objekts selbst im Absoluten\grqq{}. Die eine Unendlichkeit verschiebt das Ganze immer weiter hinaus, die andere findet das Ganze in jedem Glied. Der Satz \glqq jeder Staub eine Organisation\grqq{} ist der Naturphilosophie entnommen und sagt genau dies: Das Kleinste ist nicht bloße Materie, sondern selbst ein Sich-Organisierendes, also \SO{} in der Form des Objekts.

\paragraph{Zur Richtung der Zusammensetzung.}
Das Wort Zusammensetzung ist in einer Hinsicht zu korrigieren. Hegel sagt nicht, das Absolute bestehe aus Subjekt und Objekt wie ein Ganzes aus Stücken. Er sagt, das Absolute \glqq stellt sich in jedem der beiden dar\grqq{} und \glqq findet sich vollständig nur in beiden zusammen\grqq{}; es \glqq gebiert beide und gebiert sich aus beiden\grqq{} \Werke{94}. Die Relation von Ganzem und Seiten läuft in beide Richtungen. Bestünde das Absolute aus Subjekt und Objekt als aus Teilen, dann wäre das Subjekt nur Subjekt und das Objekt nur Objekt, und ihre Entgegensetzung wäre die ideelle. Die Zusammensetzung ist gerade dadurch eine des Absoluten, dass jeder Teil das Ganze ist. Das ist der Grund, warum die Formel (R) zu (A), (S) und (O) hinzugehört und nicht eine Ergänzung ist.

\paragraph{Zum Schlusssatz: Nur so ist Ich~=~Ich das Absolute.}
\glqq Nur indem das Objekt selbst ein \SO{} ist, ist Ich~=~Ich das Absolute.\grqq{} Der Satz kehrt Fichtes Verhältnis um. Das Selbstbewusstsein ist nicht dadurch absolut, dass es das Objekt ausschließt und sich selbst setzt, sondern dadurch, dass das Objekt seinerseits ein Sich-selbst-Setzendes ist, das Ich~=~Ich in der Form des Objekts. Erst dann kann das Ich sich im Objekt finden, und erst dann hört Ich~=~Ich auf, ein Sollen zu sein. Die Zusammensetzung ist damit die Bedingung dafür, dass das Prinzip Fichtes wahr wird, nicht seine Widerlegung.

\section{Zweiter Sachverhalt: Der Weg zur Identität der Identität und der Nichtidentität}

\subsection{Das Problem des Grundsatzes}

Das Kapitel über das Prinzip einer Philosophie in der Form eines absoluten Grundsatzes geht von der Forderung aus, dass die Philosophie als System ihr Prinzip in einem Satz aussprechen soll. Ein Satz ist aber eine Form des Verstandes. Die Vernunft, die sich für die Reflexion setzen will, kann in einem Satz die Synthese nicht ausdrücken; sie muss trennen und \glqq die Synthese und die Antithese getrennt, in zwei Sätzen, in einem die Identität, im andern die Entzweiung ausdrücken\grqq{} \Werke{37}.

\subsection{$A = A$}

Der erste Satz ist der Satz der Identität. In $A = A$ \glqq wird reflektiert auf das Bezogensein, und dies Beziehen, dies Einssein, die Gleichheit ist in dieser reinen Identität enthalten; es wird von aller Ungleichheit abstrahiert\grqq{} \Werke{37}. Dann folgt die Bemerkung, die für die spätere Formel alles entscheidet: $A = A$, \glqq Ausdruck des absoluten Denkens oder der Vernunft\grqq{}, hat \glqq für die formale, in verständigen Sätzen sprechende Reflexion nur die Bedeutung der Verstandesidentität, der reinen Einheit\grqq{}, einer solchen, \glqq worin von der Entgegensetzung abstrahiert ist\grqq{} \Werke{37}. Dasselbe Zeichen hat also zwei Bedeutungen. Als Ausdruck der Vernunft meint es die absolute Identität; für die Reflexion meint es die abstrakte Einheit, die den Gegensatz draußen lässt. Die Vernunft findet sich in dieser Einseitigkeit nicht ausgedrückt.

\subsection{$A = B$}

Darum postuliert die Vernunft \glqq auch das Setzen desjenigen, wovon in der reinen Gleichheit abstrahiert wurde, das Setzen des Entgegengesetzten, der Ungleichheit; das eine $A$ ist Subjekt, das andere Objekt, und der Ausdruck für ihre Differenz ist $A$ nicht $= A$, oder $A = B$\grqq{} \Werke{37}. Dieser Satz widerspricht dem ersten geradezu; in ihm ist die Nichtidentität gesetzt, \glqq die reine Form des Nichtdenkens\grqq{}, wie der erste die Form des reinen Denkens setzt. Beide sind \glqq ein Anderes als das absolute Denken, die Vernunft\grqq{}. Die Nichtidentität kann nur gesetzt werden, weil auch das Nichtdenken gedacht wird; das Beziehen, das Gleichheitszeichen, ist auch im zweiten Satz, aber nur subjektiv \Werke{38}.

Dann die wechselseitige Bedingtheit. \glqq Dieser zweite Satz ist so unbedingt als der erste und insofern Bedingung des ersten, so wie der erste Bedingung des zweiten Satzes ist\grqq{} \Werke{38}. Der erste besteht nur durch die Abstraktion von der Ungleichheit, die der zweite enthält; der zweite braucht, um ein Satz zu sein, das Beziehen, das der erste ausspricht. Der zweite Satz ist \glqq unter der subalternen Form des Satzes des Grundes ausgesprochen worden\grqq{} und noch tiefer herabgezogen, indem man ihn zum Satz der Kausalität gemacht hat: $A$ hat einen Grund, heißt, dem $A$ kommt ein Sein zu, das nicht ein Sein des $A$ ist, also $A$ nicht $= A$, $A = B$ \Werke{38}.

\subsection{Antinomie}

\glqq Beide Sätze sind Sätze des Widerspruchs, nur im verkehrten Sinne. Der erste, der der Identität, sagt aus, daß der Widerspruch $= 0$ ist; der zweite, insofern er auf den ersten bezogen wird, daß der Widerspruch ebenso notwendig ist als der Nichtwiderspruch\grqq{} \Werke{38~f.}. Als Sätze sind beide für sich von gleicher Potenz. Wird der zweite so ausgesprochen, dass der erste zugleich auf ihn bezogen ist, so ist er \glqq der höchstmögliche Ausdruck der Vernunft durch den Verstand; diese Beziehung beider ist der Ausdruck der Antinomie\grqq{}. Und als Antinomie, \glqq als Ausdruck der absoluten Identität\grqq{}, ist es gleichgültig, ob $A = B$ oder $A = A$ gesetzt wird: $A = A$, als Beziehung beider Sätze genommen, \glqq enthält die Differenz des $A$ als Subjekts und $A$ als Objekts zugleich mit der Identität, so wie $A = B$ die Identität des $A$ und $B$ mit der Differenz beider\grqq{} \Werke{39}. Der Verstand, der im Satz des Grundes die Antinomie nicht erkennt, ist nicht zur Vernunft gediehen; für ihn ist $A = B$ nur eine Wiederholung des $A$.

Wenn man nur auf das Formelle der Spekulation reflektiert, ist die Antinomie, \glqq der sich selbst aufhebende Widerspruch\grqq{}, der höchste formale Ausdruck des Wissens \Werke{39}. Aber die Antinomie ist nur die Seite der Reflexion. Von dieser Seite angesehen, \glqq erscheint die absolute Identität in Synthesen Entgegengesetzter, also in Antinomien\grqq{} \Werke{41}. Die andere Seite ist die transzendentale Anschauung, in der die Entgegengesetzten nicht bloß als sich aufhebend gedacht, sondern als eins angeschaut werden. Spekulation ist die Einheit beider Seiten.

\subsection{Fichtes Entscheidung}

An dieser Stelle lässt sich Fichtes System als ein bestimmter Umgang mit den beiden Sätzen beschreiben. \glqq $A = A$ und $A = B$ bleiben beide unbedingt; es soll nur $A = A$ gelten; d.\,h. aber, ihre Identität ist nicht in ihrer wahren Synthese, die kein bloßes Sollen ist, dargestellt. So ist im Fichteschen System Ich~=~Ich das Absolute\grqq{} \Werke{50}. Die Totalität der Vernunft führt den zweiten Satz herbei, der ein Nicht-Ich setzt; die Synthese wird postuliert, aber in ihr bleibt die Entgegensetzung: \glqq der eine Satz soll bestehen, der eine höher an Rang sein als der andere\grqq{}. Die Antinomie wird nicht aufgehoben, sondern als Streben festgehalten. Das Sollen ist die Gestalt, in der die nicht vollzogene Identität der beiden Sätze erscheint.

\begin{figure}[!htb]
\centering
\resizebox{\linewidth}{!}{%
\begin{tikzpicture}[
  every node/.style={align=center, font=\small},
  satz/.style={draw, rounded corners=2pt, inner sep=5pt, minimum width=4cm},
  edge/.style={-{Latex[length=2mm]}, thick}]
\node[satz] (AA) at (-4.6,0) {$A = A$\\ Satz der Identität\\ \footnotesize Widerspruch $= 0$\\ \footnotesize für die Reflexion: Verstandesidentität};
\node[satz] (AB) at (4.6,0) {$A = B$ \; ($A \neq A$)\\ Satz des Grundes, Nichtidentität\\ \footnotesize Widerspruch so notwendig\\ \footnotesize wie Nichtwiderspruch};
\draw[edge, <->, >={Latex[length=2mm]}] (AA) -- node[above=0.5mm, font=\scriptsize] {jeder unbedingt,\\ jeder Bedingung\\ des anderen} (AB);
\node[draw, rounded corners=2pt, inner sep=5pt] (ANT) at (0,-2.4) {\textbf{Antinomie}\\ \footnotesize der sich selbst aufhebende Widerspruch\\ \footnotesize höchster Ausdruck der Vernunft durch den Verstand};
\draw[edge] (AA.south) -- (ANT.north west);
\draw[edge] (AB.south) -- (ANT.north east);
\node[draw, rounded corners=2pt, inner sep=5pt, fill=gray!12] (ID) at (0,-4.9) {\textbf{Das Absolute: Identität der Identität und der Nichtidentität}\\ \footnotesize \glqq Entgegensetzen und Einssein ist zugleich in ihm\grqq{}\\ \footnotesize Antinomie (Reflexion) und transzendentale Anschauung (Vernunft) in einem};
\draw[edge] (ANT) -- node[right, font=\scriptsize] {nicht bloß gedacht,\\ sondern angeschaut} (ID);
\node[draw, dashed, inner sep=3pt, font=\scriptsize] (F) at (-4.6,1.9) {Fichte: \glqq es soll nur $A = A$ gelten\grqq{}\\ Ich = Ich wird zu \glqq Ich soll gleich Ich sein\grqq{}};
\draw[dashed, gray] (F) -- (AA);
\end{tikzpicture}}
\caption{Der Weg zur Formel nach W~2, S.~37--39, 41, 50 und 96. Die beiden Grundsätze bedingen einander; ihre Beziehung ist die Antinomie; das Absolute ist die Identität, in der beide Sätze zugleich gelten, statt dass einer gelten soll.}
\label{fig:antinomie}
\end{figure}

\subsection{Die Formel}

In der Vergleichung kehrt die Zweiheit der Sätze als Zweiheit von Trennung und Identität wieder. \glqq Die Philosophie muß dem Trennen in Subjekt und Objekt sein Recht widerfahren lassen; aber indem sie es gleich absolut setzt mit der der Trennung entgegengesetzten Identität, hat sie es nur bedingt gesetzt, so wie eine solche Identität -- die durch Vernichten der Entgegengesetzten bedingt ist -- auch nur relativ ist\grqq{} \Werke{96}. Dann der Satz:

\begin{quote}
\glqq Das Absolute selbst aber ist darum die Identität der Identität und der Nichtidentität; Entgegensetzen und Einssein ist zugleich in ihm.\grqq{} \Werke{96}
\end{quote}

Das \glqq darum\grqq{} verweist auf die Bedingtheit: Trennung, die der Identität gleich absolut gesetzt wird, ist bedingt; Identität, die durch Vernichten der Entgegengesetzten zustande kommt, ist relativ. Weder die Trennung für sich noch die Identität für sich ist das Absolute. Das Absolute ist das, worin beide sind, und zwar nicht nacheinander, sondern zugleich. Zwei Seiten später wird dieselbe Sache von der reellen Entgegensetzung her gesagt: Sie ist Werk der Vernunft, \glqq welche die Entgegengesetzten nicht bloß in der Form des Erkennens, sondern auch in der Form des Seins, Identität und Nichtidentität identisch setzt\grqq{} \Werke{98}.

\subsection{Die Formel, kommentiert}

\paragraph{Zwei Bedeutungen des einen Wortes.}
Die Formel enthält das Wort Identität zweimal, und die beiden Vorkommen bedeuten nicht dasselbe. Das ist keine Nachlässigkeit, sondern die Anwendung dessen, was das Grundsatzkapitel über $A = A$ gesagt hat: Dasselbe Zeichen bedeutet für die Reflexion die Verstandesidentität und als Ausdruck der Vernunft die absolute Identität \Werke{37}. Die innere Identität der Formel ist $A = A$, die relative Identität, die von der Entgegensetzung abstrahiert. Die Nichtidentität ist $A = B$. Die äußere Identität ist die der Vernunft, die beide zusammenhält. Die Formel schreibt beide Bedeutungen in einen Ausdruck und ordnet sie: Die Vernunftidentität ist die Identität, in der die Verstandesidentität und ihre Negation beide stehen.

\paragraph{Warum nicht bloß Identität.}
Man könnte meinen, das Absolute sei einfach Identität. Aber die reine Identität ist durch die Abstraktion von der Ungleichheit bedingt; der erste Satz ist durch den zweiten bedingt \Werke{38}. Eine Identität, die nur dadurch besteht, dass sie das Entgegengesetzte vernichtet, ist relativ \Werke{96}. Sie ist eine Seite und darum nicht das Absolute. Das gilt auch dann, wenn diese eine Seite, wie bei Fichte, das Selbstbewusstsein ist.

\paragraph{Warum nicht bloß Identität und Nichtidentität.}
Das Nebeneinander beider Sätze ist die Antinomie, der höchste Ausdruck der Vernunft durch den Verstand. Aber als bloßes Nebeneinander sind es zwei Sätze, die einander bedingen, und die Frage, welcher gelten soll, bleibt offen. Fichte beantwortet sie zugunsten von $A = A$ und verwandelt damit die Identität in ein Sollen \Werke{50}. Die äußere Identität der Formel sagt, dass das Zugleich beider kein Sollen ist, sondern ist: \glqq Entgegensetzen und Einssein ist zugleich in ihm\grqq{}. Das Und der Formel ist die Antinomie; der Genitiv, Identität \emph{der}, ist die Anschauung, die das Und als Eines fasst.

\paragraph{Die Selbstanwendung.}
Die Formel wendet die Identität auf sich selbst und auf ihre Negation an. Das ist der Grund, warum sie nicht vom Satz des Widerspruchs getroffen wird. Die beiden Sätze sind \glqq Sätze des Widerspruchs, nur im verkehrten Sinne\grqq{} \Werke{38}; der eine sagt, der Widerspruch sei null, der andere, er sei ebenso notwendig wie der Nichtwiderspruch. Wer beide in einem denkt, denkt den \glqq sich selbst aufhebenden Widerspruch\grqq{} \Werke{39}. Die Formel benennt das Ganze, in dem dieser Widerspruch sich aufhebt, ohne dass die Entgegengesetzten verschwinden. Das ist der Unterschied zur Vernichtung: Die Getrennten \glqq sollen als Getrennte bleiben und diesen Charakter nicht verlieren, insofern sie im Absoluten oder das Absolute in ihnen gesetzt ist\grqq{} \Werke{96}.

\paragraph{Die Formel und die Zusammensetzung.}
Damit wird der Zusammenhang der beiden Sachverhalte sichtbar. Die Formel ist die Satzform der Struktur, die der erste Sachverhalt beschreibt.

\begin{center}
\renewcommand{\arraystretch}{1.3}
\begin{tabular}{@{}>{\raggedright\arraybackslash}p{3.1cm} >{\raggedright\arraybackslash}p{4.6cm} >{\raggedright\arraybackslash}p{6.6cm}@{}}
\toprule
Moment der Formel & Struktur der Zusammensetzung & Wortlaut \\
\midrule
innere Identität & jede Seite ist für sich \SO{} & \glqq das Subjekt subjektives \SO{}, das Objekt objektives \SO{}\grqq{} (S.~96) \\
Nichtidentität & die reelle Entgegensetzung der beiden Formen & \glqq zugleich beide als Getrennte bestehen\grqq{} (S.~96) \\
äußere Identität & der Indifferenzpunkt, gedoppelt als Identität und Totalität & \glqq beide gebiert und sich aus beiden gebiert\grqq{} (S.~94, 111) \\
\bottomrule
\end{tabular}
\end{center}

Die Verdopplung des Mittelpunkts in Identität und Totalität \Werke{111} ist dieselbe Verdopplung wie die des Wortes Identität in der Formel: Identität als das, was in jedem Pol liegt, und Identität als das Ganze, das beide Pole sind. Und die Rekursion des ersten Sachverhalts findet ihr Gegenstück darin, dass die innere Identität der Formel, wenn sie auf einen Teil angewendet wird, selbst wieder die Identität der Identität und der Nichtidentität ist, denn jeder Teil ist im Absoluten. Hegel iteriert die Formel nicht ausdrücklich; er iteriert das \SO{}. Dass beides dieselbe Sache ist, ergibt sich aus dem Satz, der beide Sachverhalte verbindet: Nur weil das Objekt selbst \SO{} ist, sind Identität und Nichtidentität \glqq auch in der Form des Seins\grqq{} identisch gesetzt und nicht bloß in der Form des Erkennens \Werke{98}.

\section{Ergebnis}

Die Differenzschrift entwickelt beide Sachverhalte aus einem einzigen Mangel des Fichteschen Systems: Dort ist die Identität nur unter der Form des Subjekts gesetzt, und darum wird sie zum Sollen. Die Abhilfe ist die doppelte Setzung. Subjekt und Objekt sind beide \SO{}, das Absolute stellt sich in jedem dar und ist nur beide zusammen; die Teilung geht ins Unendliche, ohne die Identität zu verlieren. In Satzform lautet dieselbe Sache: Das Absolute ist die Identität der Identität und der Nichtidentität. Die innere Identität ist die relative, die jede Seite für sich hat; die Nichtidentität ist ihre reelle Entgegensetzung; die äußere Identität ist der Indifferenzpunkt, der beide gebiert und sich aus beiden gebiert. Die Formel ist damit weder Tautologie noch Widerspruch, sondern die Angabe, auf welcher Ebene die Antinomie der beiden Grundsätze aufhört, ein Sollen zu sein.

\clearpage
\section*{Apparat}

\subsection*{1. Statustafel}
\begin{center}
\renewcommand{\arraystretch}{1.3}
\begin{tabular}{@{}p{8.6cm}p{2.4cm}p{3.6cm}@{}}
\toprule
Aussage & Status & Beleg \\
\midrule
Subjekt und Objekt werden beide als \SO{} gesetzt; das Absolute stellt sich in jedem dar und ist nur beide zusammen. & belegt & W~2, S.~94 \\
Die Zweiheit kehrt in jeder Seite und in jedem Teil wieder (Teilung ins Unendliche). & belegt & W~2, S.~96~f. \\
Der Unterschied der beiden Seiten ist der des Substantialitätsverhältnisses (Formindex), nicht des Gehalts. & belegt & W~2, S.~101 \\
Der Mittelpunkt ist gedoppelt: Identität und Totalität. & belegt & W~2, S.~111 \\
Die Formel steht in der Vergleichung; ihre Bausteine stammen aus dem Grundsatzkapitel. & belegt & W~2, S.~37--39, 96 \\
Das Wort Identität hat in der Formel zwei Bedeutungen (Verstandes- und Vernunftidentität). & Deutung, gestützt & W~2, S.~37 \\
Die Formel ist die Satzform der Zusammensetzung (Zuordnung der drei Momente). & Deutung & W~2, S.~94, 96, 98, 111 \\
\glqq Zusammensetzung\grqq{} ist als wechselseitige Erzeugung, nicht als Bestehen aus Teilen zu lesen. & Deutung, gestützt & W~2, S.~94 \\
Die Formel selbst wird rekursiv iteriert. & offen; Hegel iteriert das \SO{}, nicht die Formel & W~2, S.~96~f. \\
Seitenangabe des Zitats zur Bedingtheit des Ich durch das Objekt. & unscharf (S.~68--73) & Paginierung der Vorlage lückenhaft \\
\bottomrule
\end{tabular}
\end{center}

\subsection*{2. Abbildungen}
\begin{itemize}
\item Abbildung~\ref{fig:baum}: Die Zusammensetzung als Baum mit doppelter Erzeugungsrichtung und Teilung ins Unendliche.
\item Abbildung~\ref{fig:pole}: Transzendentalphilosophie und Naturphilosophie als Pole der Indifferenz mit gedoppeltem Mittelpunkt.
\item Abbildung~\ref{fig:antinomie}: Die beiden Grundsätze, ihre Antinomie und die Formel; Fichtes Entscheidung für $A = A$.
\end{itemize}

\subsection*{3. Sachen}
\begin{itemize}
\item \SO{}
\item subjektives \SO{}
\item objektives \SO{}
\item Indifferenzpunkt
\item Identität und Totalität (gedoppelter Mittelpunkt)
\item Entzweiung
\item ideelle und reelle Entgegensetzung
\item Substantialitätsverhältnis
\item Teilung ins Unendliche
\item Progress ins Unendliche (schlechte Unendlichkeit)
\item Satz der Identität ($A = A$)
\item Satz des Grundes ($A = B$)
\item Antinomie
\item transzendentale Anschauung
\item Verstandesidentität und Vernunftidentität
\item Sollen
\item Transzendentalphilosophie
\item Naturphilosophie
\item Form und Materie
\end{itemize}

\subsection*{4. Personen}
\begin{itemize}
\item Georg Wilhelm Friedrich Hegel
\item Johann Gottlieb Fichte
\item Friedrich Wilhelm Joseph Schelling
\item Immanuel Kant
\item Karl Leonhard Reinhold
\item Platon
\end{itemize}

\subsection*{5. Quellen}
\begin{itemize}
\item Hegel, G.\,W.\,F.: Differenz des Fichteschen und Schellingschen Systems der Philosophie (1801). In: Werke, Band~2: Jenaer Schriften 1801--1807. Frankfurt am Main: Suhrkamp 1986, S.~9--138. Zitiert als W~2.
\item Fichte, J.\,G.: Grundlage der gesamten Wissenschaftslehre. Sämtliche Werke, Band~I, S.~253. Zitiert nach Hegels Anmerkung, W~2, S.~99.
\item Platon: Timaios 31--32. Zitiert nach Hegels Anmerkung, W~2, S.~98.
\end{itemize}

\subsection*{Dokumentdaten}
\begin{itemize}
\item Titel: Subjekt-Objekt und die Identität der Identität und der Nichtidentität. Die Zusammensetzung des Absoluten, des Subjektiven und des Objektiven in Hegels Differenzschrift (1801)
\item Datei: subjekt-objekt\_identitaet\_der\_identitaet\_v1.tex
\item Version: v1
\item Datum: 2026-09-25
\item Grundlage: Hegel, Werke Band~2 (Suhrkamp 1986), Textdatei der Vorlage; Seitenzahlen nach eingedruckter Paginierung
\item Status: Erstfassung; Kanon nach Pauschale, solange die Qualität hält
\end{itemize}

\end{document}
